\documentclass[11pt]{article}
\usepackage{preamble}
\renewcommand{\answer}[1]{}

\begin{document}

\ladderheading{AP Statistics \textperiodcentered\ Unit 2 \textperiodcentered\ Topic 2.5 \textperiodcentered\ B: 2026-10-23 \textperiodcentered\ E: 2026-11-09}{Count the Overlap Once}

\focusbanner{TODAY'S PROBLEM}

Suppose 120 students answer two yes-or-no questions: ``Do you like pizza?'' and ``Do you like tacos?'' The table summarizes their answers. Let $P$ be the event that a randomly selected student likes pizza and $T$ the event that the student likes tacos.\par\smallskip \textbf{Jamie:} ``$P(P\text{ or }T)=80/120+60/120=140/120$. Because a probability cannot exceed 1, this must mean everyone likes at least one; I would report the answer as 1.''\par \textbf{Noor:} ``Some students may have been counted in both totals, so we need the intersection before deciding.''\par
\begin{center}
\textbf{Student food preferences (counts)}\par\smallskip
\begin{tabular}{|l|c|c|c|}
\hline
\textbf{Preference} & \textbf{Pizza} & \textbf{Not pizza} & \textbf{Total} \\ \hline
\textbf{Tacos} & 35 & 25 & 60 \\ \hline
\textbf{Not tacos} & 45 & 15 & 60 \\ \hline
\textbf{Total} & 80 & 40 & 120 \\ \hline
\end{tabular}
\end{center}
\begin{firsttakebox}{FIRST TAKE --- before the video (5 min)}
Does the table support reporting everyone likes pizza or tacos? Choose a claim and cite one count.
\workspace{0.9in}
\end{firsttakebox}

\begin{callout}{\IconBook\ INTERSECTION, UNION, AND DISJOINT EVENTS (keep on the page)}{calloutgreen}
$A\cap B$ is the \textbf{intersection}: outcomes in both $A$ and $B$.
Events are \textbf{mutually exclusive} (disjoint) exactly when they cannot happen together, so $P(A\cap B)=0$.
$A\cup B$ is the \textbf{union}: outcomes in $A$, in $B$, or in both. Count every outcome once:
$P(A\cup B)=P(A)+P(B)-P(A\cap B)$. The subtraction removes the second count of the overlap.
Every probability must be between 0 and 1; a result outside that range signals an error to diagnose, not a value to cap.
\end{callout}

\newpage
\textbf{Finish after the video:}\par\smallskip

\dokbadge{1}~\textbf{(a)} State the number of students in $P\cap T$ and write $P(P\cap T)$ as a fraction.\par
\workspace{0.7in}

\dokbadge{2}~\textbf{(b)} Compute $P(P\cup T)$ so every student in the union is counted once. Also compute the probability that a student likes neither food.\par
\workspace{1.3in}

\dokbadge{3}~\textbf{(c)} Decide whether $P$ and $T$ are mutually exclusive and which student's reasoning is warranted. Cite the 35-student intersection and your union probability, explain what Jamie's calculation would mislead a reader into believing, and name a pair of events from this context that really are mutually exclusive.\par
\workspace{2.0in}

\begin{sentenceframebox}\raggedright\textbf{Frame:} The events are \blank[1.0in] mutually exclusive because the intersection contains \blank[0.6in] students, so \blank[2.0in].\end{sentenceframebox}

\begin{sentenceframebox}\raggedright\textbf{Frame:} The sum \blank[1.0in] counts \blank[2.5in] twice; reporting it would make a reader think \blank[2.4in].\end{sentenceframebox}

\begin{summaryexitbox}{EXIT --- turn this in}
One thing the video changed about my first take: \hrulefill\par
\workspace{0.4in}
\end{summaryexitbox}

\end{document}
